\documentclass[12pt]{article}
\usepackage[utf8]{inputenc}
\usepackage[T1]{fontenc}
\usepackage[margin=1in]{geometry}
\usepackage{setspace}
\usepackage{xcolor}
\usepackage[colorlinks=true, urlcolor=blue!60!black, linkcolor=blue!60!black]{hyperref}
\setlength{\parindent}{2em}
\setlength{\parskip}{1.25em}
\renewcommand{\baselinestretch}{1.0}
\begin{document}
\begin{center}
{\LARGE \textit{New York State Rifle \& Pistol Association, Inc. v. Bruen}, 597 U.S. 1 (2022)\par}
\vspace{0.8em}
{\large Case Brief --- Four-Page Edition\par}
\end{center}
\vspace{1em}

\section{Memory Jogger}

The Second and Fourteenth Amendments protect an ordinary, law-abiding citizen's right to carry a handgun in public for self-defense, and regulations of protected conduct are valid only if the government proves them consistent with the Nation's historical tradition of firearm regulation, not under means-end scrutiny (\href{https://www.law.cornell.edu/supremecourt/text/20-843}{597 U.S. 1}).

\section{Detailed Case Facts}

New York has regulated public carry of handguns since the 1905 ban on unlicensed concealed carry and the 1911 Sullivan Law; a 1913 amendment required ``good moral character'' and ``proper cause'' for a concealed-carry license. Because New York banned open carry, an unrestricted concealed-carry license was the only lawful way for most residents to carry a handgun for self-defense. State courts read ``proper cause'' to require a special need for self-protection distinguishable from that of the general community, which excluded generalized safety concerns. Forty-three states used objective shall-issue regimes; six states and the District of Columbia used discretionary may-issue regimes.

Brandon Koch and Robert Nash, law-abiding Rensselaer County residents and members of the New York State Rifle \& Pistol Association, held licenses restricted to hunting and target shooting. A licensing officer refused to lift the restrictions for Nash, who cited neighborhood robberies, and expanded Koch's only to carrying to and from work.

\section{Procedural History}

Petitioners sued under \S~1983. Bound by \textit{Kachalsky v.\ County of Westchester}, which had upheld the proper-cause requirement under intermediate scrutiny, the Northern District of New York dismissed, and the Second Circuit summarily affirmed. The Supreme Court reversed and remanded on June 23, 2022.

\section{Judicial Votes}

Reversed and remanded, \textbf{6--3}.
\begin{itemize}
\item \textbf{Majority:} Justice Thomas, joined by Chief Justice Roberts and Justices Alito, Gorsuch, Kavanaugh, and Barrett.
\item \textbf{Concurring:} Justice Alito; Justice Kavanaugh, joined by Chief Justice Roberts; Justice Barrett.
\item \textbf{Dissenting:} Justice Breyer, joined by Justices Sotomayor and Kagan.
\end{itemize}

\section{Holding}

The Second and Fourteenth Amendments protect the right to carry a handgun for self-defense outside the home. Rejecting the courts of appeals' two-step framework, the Court held that when the plain text covers conduct, the government must show that its regulation is consistent with historical tradition; for unprecedented concerns, courts reason by analogy, comparing how and why regulations burden armed self-defense. New York's proper-cause requirement is unconstitutional. Footnote 9's approval of objective shall-issue regimes, the remarks on sensitive places, and the open 1791-versus-1868 question are dicta or reservations.

\section{Analysis of Opinions}

\textbf{Majority (Thomas).} The two-step framework was ``one step too many'': interest balancing contradicts \textit{Heller} and \textit{McDonald}, because the Second Amendment is itself the product of balancing by the people. For persistent problems, the lack of a similar historical regulation is evidence of unconstitutionality; for new ones, courts seek a representative analogue. Koch and Nash are among ``the people,'' handguns are in common use, and ``bear'' includes public carry. New York's history failed: the Statute of Northampton came to bar only carrying to terrify the public; colonial laws targeted terrorizing conduct; antebellum concealed-carry bans were upheld only where open carry remained, and surety laws presumed a right to carry; Reconstruction confirmed the right; and late Western territorial bans were outliers.

\textbf{Alito (concurring).} The holding is narrow and leaves \textit{Heller} intact; the dissent's statistics do not show that the law deterred criminals.

\textbf{Kavanaugh (concurring).} States may still require licenses through objective shall-issue regimes with fingerprinting, background checks, and training, and longstanding prohibitions remain valid.

\textbf{Barrett (concurring).} The Court left open how post-ratification practice informs original meaning and whether 1791 or 1868 controls.

\textbf{Breyer (dissenting).} Citing 45,222 firearm deaths in 2020, Breyer argued that legislatures need leeway. The Court decided the case on the pleadings without a record, replaced a unanimous framework with a rigid history-only test, and left judges who are not historians to ``pick their friends out of history's crowd.''

\section{Examples: Future Applications}

\begin{enumerate}
\item \textbf{Same-Side (Invalid): A Discretionary Character Test.} A state replaces proper cause with an open-ended ``good moral character'' standard lacking objective criteria. It replicates the unchanneled discretion \textit{Bruen} struck down.
\item \textbf{Same-Side (Invalid): Sweeping Sensitive Places.} A state bans carry in all businesses open to the public, downtown sidewalks, and transit hubs. Like declaring Manhattan a sensitive place, it lacks a historical analogue.
\item \textbf{Opposite-Side (Valid): Objective Shall-Issue Licensing.} A state requires age 21, background checks, fingerprints, a modest fee, and a safety course, then must issue within 30 days. Footnote 9 and Kavanaugh's concurrence approve such regimes.
\item \textbf{Opposite-Side (Valid): Disarming Credible Threats.} A ban on possession by persons under domestic-violence orders with a credible-threat finding is relevantly similar to surety and going-armed laws, as \textit{United States v.\ Rahimi} held.
\item \textbf{Fence-Sitter: The Default Private-Property Rule.} A state bars carry on private property open to the public unless the owner consents. Proponents cite historical trespass laws; challengers say those covered enclosed lands, not commercial venues.
\end{enumerate}

\section{Critique}

\textbf{Originalist perspectives.} Mark W. Smith praised \textit{Bruen} for enforcing \textit{Heller}'s method. Yet William Baude and Robert Leider argued that the search for statutory analogues overlooks the Founding-era general-law framework, and Randy Barnett and Lawrence Solum questioned the tension between original public meaning and a free-standing history-and-tradition test.

\textbf{Progressive perspectives.} Joseph Blocher, Eric Ruben, and Jacob Charles criticized ``originalism-by-analogy'' for confusing lower courts about levels of generality, the meaning of historical silence, and judges' competence to do history. Feminist scholars warned that tying law to eras when domestic violence was unregulated threatens vulnerable people, and \textit{Rahimi} soon clarified that \textit{Bruen} seeks consistency with underlying principles, not a historical twin.

\section{Key Quotations}

\begin{quote}
``When the Second Amendment's plain text covers an individual's conduct, the Constitution presumptively protects that conduct. The government must then justify its regulation by demonstrating that it is consistent with the Nation's historical tradition of firearm regulation. Only then may a court conclude that the individual's conduct falls outside the Second Amendment's `unqualified command.' ''
\end{quote}
\noindent --- Justice Thomas, Opinion of the Court, 597 U.S. 1.

\begin{quote}
``To be clear, analogical reasoning under the Second Amendment is neither a regulatory straightjacket nor a regulatory blank check. On the one hand, courts should not `uphold every modern law that remotely resembles a historical analogue,' because doing so `risk[s] endorsing outliers that our ancestors would never have accepted.' On the other hand, analogical reasoning requires only that the government identify a well-established and representative historical analogue, not a historical twin. So even if a modern-day regulation is not a dead ringer for historical precursors, it still may be analogous enough to pass constitutional muster.''
\end{quote}
\noindent --- Justice Thomas, Opinion of the Court, 597 U.S. 1.

\begin{quote}
``Going forward, therefore, the 43 States that employ objective shall-issue licensing regimes for carrying handguns for self-defense may continue to do so. Likewise, the 6 States including New York potentially affected by today's decision may continue to require licenses for carrying handguns for self-defense so long as those States employ objective licensing requirements like those used by the 43 shall-issue States.''
\end{quote}
\noindent --- Justice Kavanaugh, concurring, joined by Chief Justice Roberts, 597 U.S. 1.

\end{document}
